%!TEX program = xelatex

\documentclass[compress]{beamer}
%--------------------------------------------------------------------------
% Common packages
%--------------------------------------------------------------------------
\usepackage[english]{babel}
\usepackage{pgfpages} % required for notes on second screen
\usepackage{graphicx}

\usepackage{multicol}

\usepackage{tabularx,ragged2e}
\usepackage{booktabs}


%--------------------------------------------------------------------------
% Load theme
%--------------------------------------------------------------------------
\usetheme{hri}

\usepackage{dtklogos} % must be loaded after theme
\usepackage{tikz}
\usetikzlibrary{mindmap,backgrounds,positioning}

\graphicspath{{figs/}}
\setbeamercolor{refToContribCol}{bg=hriSec1Comp,fg=white}
\newcommand{\refToContrib}[1]{%
    \begin{beamercolorbox}[wd=\linewidth,ht=2ex,dp=0.7ex]{refToContribCol}%
    \hspace{0.5em}$\hookrightarrow$ #1%
    \end{beamercolorbox}%
}%

%--------------------------------------------------------------------------
% General presentation settings
%--------------------------------------------------------------------------
\title{Social Interaction Scenarios}
\subtitle{Q6 -- Can you devise a social interaction scenario that current
cognitive architectures would likely fail, and why?}
\date{}
\author{}
\institute{\includegraphics[height=15mm]{plymouth-logo}}
%\\Centre for Robotics \& Neural
%Systems\\{\Medium Plymouth University}}

%--------------------------------------------------------------------------
% Notes settings
%--------------------------------------------------------------------------
%\setbeameroption{show notes on second screen}

\begin{document}
%--------------------------------------------------------------------------
% Titlepage
%--------------------------------------------------------------------------

\maketitle

%\imageframe[Children playing with\\the Ranger robot]{photo-fullscreen.jpg}
%\imageframe{photo-fullscreen.jpg}

\begin{frame}{Cognitive Architecture?}
\centering
        \resizebox{!}{0.7\paperheight}{%
            \begin{tikzpicture}[
                    >=latex,
                every edge/.style={<-, draw, very thick}]
        

            \path[small mindmap, 
                level 1 concept/.append style={sibling angle=360/5}, 
                level 2 concept/.append style={sibling angle=120}, 
            concept color=hriWarmGreyLight,text=hriWarmGreyDark]
            node[concept] {Cognition}
            [clockwise from=-30]
            child[concept color=hriSec1,text=white] { node[concept] (percept) {Perception}
                [clockwise from=40]
                child[concept color=hriSec2Dark,text=white] { node[concept]{Attention} }
                child[concept color=hriSec2CompDark,text=white] { node[concept] (dialog) {Communication} }
            }
            child[concept color=hriSec2Comp,text=white] { node[concept] (krr) {Knowledge Representation} 
                [clockwise from=-30]
                child[concept color=hriSec1CompDark,text=white] { node[concept] (memory) {Memory} }
                child[concept color=hriSec3CompDark,text=white] { node[concept] (tom) {Theory of Mind} }
            }
            child[concept color=hriSec2,text=white] { node[concept] (action) {Action Execution} 
                [counterclockwise from=110]
                child[concept color=hriSec3,text=white] { node[concept] (plan) {Task Planning} }
                child[concept color=hriSec2CompDark,text=white] { node[concept] (comm) {Communication} }
            }
            child[concept color=hriSec3Comp,text=white] { node[concept] (learning) {Learning} 
                [counterclockwise from=130]
                child[concept color=hriSec1CompDark,text=white] { node[concept] (adapt) {Adaptation} }
            }
            child[concept color=hriSec3CompDark,text=white] { node[concept] (reason) {Reasoning} 
                [clockwise from=110]
                child[concept color=hriSec3Dark,text=white] { node[concept] {Problem solving} }
                child[concept color=hriSec1Dark,text=white] { node[concept] (decision) {Decision making} } 
            };

        \onslide<2>{
        \path (decision) edge[->, bend left] (action);
        \path (percept) edge[->, bend left] (action);
        \path (krr) edge[<->, bend right] (reason);
        \path (plan) edge[<->] (reason);
        \path (percept) edge[->, bend left] (krr);
        \path (krr) edge[->, bend left] (plan);
        \path (learning) edge[->, bend left] (reason);
        \path (learning) edge[->, bend left] (krr);
        \path (percept) edge[<->, bend left] (learning);
        \path (dialog) edge[->, bend left] (tom);
        \path (tom) edge[->, bend left] (comm);
        \path (percept) edge[<->, bend left] (memory);
        }


        \end{tikzpicture}
    }
\end{frame}

\begin{frame}{}

    \centering

    \begin{exampleblock}{Cognitive Architectures Building $\neq$ Science of Integration}
        The function may be a {\Medium byproduct} of the architecture
    \end{exampleblock}
    \vspace{2em}
    \uncover<2>{
        At this stage, we should not commit to any architecture,\\nor internal mechanism
    }


\end{frame}


\begin{frame}{No commitments to internal mechanisms}
    \centering
        \resizebox{!}{0.7\paperheight}{%
            \begin{tikzpicture}[
                    >=latex,
                every edge/.style={<-, draw, very thick}]
        
         \onslide<1>{
            \path[small mindmap, 
                level 1 concept/.append style={sibling angle=360/5}, 
                level 2 concept/.append style={sibling angle=120}, 
            concept color=hriWarmGreyLight,text=hriWarmGreyDark]
            node[concept color=white] {}
            [clockwise from=-30]
            child[concept color=hriSec1,text=white] { node[concept] (percept) {Perception}
                [clockwise from=40]
                child[concept color=hriSec2Dark,text=white] { node[concept]{Attention}; }
                child[concept color=hriSec2CompDark,text=white] { node[concept] (dialog) {Communication} ;};
            }
            child[concept color=hriSec2Comp,text=white] { node[concept] (krr) {Knowledge Representation} 
                [clockwise from=-30]
                child[concept color=hriSec1CompDark,text=white] { node[concept] (memory) {Memory}; }
                child[concept color=hriSec3CompDark,text=white] { node[concept] (tom) {Theory of Mind}; };
            }
            child[concept color=hriSec2,text=white] { node[concept] (action) {Action Execution} 
                [counterclockwise from=110]
                child[concept color=hriSec3,text=white] { node[concept] (plan) {Task Planning}; }
                child[concept color=hriSec2CompDark,text=white] { node[concept] (comm) {Communication}; };
            }
            child[concept color=hriSec3Comp,text=white] { node[concept] (learning) {Learning} 
                [counterclockwise from=130]
                child[concept color=hriSec1CompDark,text=white] { node[concept] (adapt) {Adaptation}; };
            }
            child[concept color=hriSec3CompDark,text=white] { node[concept] (reason) {Reasoning} 
                [clockwise from=110]
                child[concept color=hriSec3Dark,text=white] { node[concept] {Problem solving};}
                child[concept color=hriSec1Dark,text=white] { node[concept] (decision) {Decision making}; } ;
            };
        }
         \onslide<2>{
            \path[small mindmap, 
                level 1 concept/.append style={sibling angle=360/5}, 
                level 2 concept/.append style={sibling angle=120}, 
            concept color=hriWarmGreyLight,text=hriWarmGreyDark]
            node[concept color=white] {}
            [clockwise from=-30]
            child[concept color=hriSec1,text=white] { node[concept] (percept) {Perception}
                [clockwise from=40]
                child[concept color=hriSec2Dark,text=white] { node[concept]{Attention}; }
                child[concept color=hriSec2CompDark,text=white] { node[concept] (dialog) {Communication} ;};
            }
            child[concept color=white,text=gray] { node[concept] (krr) {Knowledge Representation} 
                [clockwise from=-30]
                child[concept color=white,text=gray] { node[concept] (memory) {Memory}; }
                child[concept color=hriSec3CompDark,text=white] { node[concept] (tom) {Theory of Mind}; };
            }
            child[concept color=hriSec2,text=white] { node[concept] (action) {Action Execution} 
                [counterclockwise from=110]
                child[concept color=white,text=gray] { node[concept] (plan) {Task Planning}; }
                child[concept color=hriSec2CompDark,text=white] { node[concept] (comm) {Communication}; };
            }
            child[concept color=white,text=gray] { node[concept] (learning) {Learning} 
                [counterclockwise from=130]
                child[concept color=hriSec1CompDark,text=white] { node[concept] (adapt) {Adaptation}; };
            }
            child[concept color=white,text=gray] { node[concept] (reason) {Reasoning} 
                [clockwise from=110]
                child[concept color=white,text=gray] { node[concept] {Problem solving};}
                child[concept color=white,text=gray] { node[concept] (decision) {Decision making}; } ;
            };
        }

        \end{tikzpicture}
    }


\end{frame}

\begin{frame}{}

    \centering

    \begin{exampleblock}{Cognitive Architectures Building $\neq$ Science of Integration}
        The function may be a {\Medium byproduct} of the architecture
    \end{exampleblock}
    \vspace{2em}
    What are the {\Medium surface} cognitive competencies\\required by {\Medium social interaction}?


\end{frame}



\begin{frame}{Functions for Social Cognition}
\centering
        \resizebox{!}{0.7\paperheight}{%
            \begin{tikzpicture}[
                    >=latex,
                every edge/.style={<-, draw, very thick}]
        

            \path[small mindmap, 
                level 1 concept/.append style={sibling angle=360/6}, 
                level 2 concept/.append style={sibling angle=60}, 
            concept color=hriWarmGreyLight,text=hriWarmGreyDark]
            node[concept] {Social\\Cognition in HRI}
            [clockwise from=30]
            child[concept color=hriSec1,text=white] { node[concept] (percept) {Perception of Human's State}
                [clockwise from=120]
                child[concept color=hriSec3Dark,text=white] { node[concept]{Emotions} }
                child[concept color=hriSec2Dark,text=white] { node[concept]{Attention} }
                child[concept color=hriSec2CompDark,text=white] { node[concept] {Inference of mental models} }
            }
            child[concept color=hriSec2Comp,text=white,grow=-30] { node[concept] (knowledge) {Social Knowledge} 
                [counterclockwise from=-120]
                child[concept color=hriSec1CompDark,text=white] { node[concept] {Social rules} }
                child[concept color=hriSec3Comp,text=black] { node[concept] {Social context} }
                child[concept color=hriSec3CompDark,text=white] { node[concept] {Common-sense} }
            }
            child[concept color=hriSec3Comp,text=black, grow=-120] { node[concept] (comm) {Communication} 
                [counterclockwise from=180]
                child[concept color=hriSec1CompDark,text=white] { node[concept] {Verbal} }
                child[concept color=hriSec1Dark,text=white] { node[concept] {Non-verbal} }
            }
            child[concept color=hriSec3,text=white,grow=180] { node[concept] (dynamics) {Interaction Dynamics} 
                [clockwise from=180]
                child[concept color=hriSec2Dark,text=white] { node[concept] {Long-term interaction} }
            }
            child[concept color=hriSec2,text=black, grow=120] { node[concept] (action) {Human-aware actions} 
                [counterclockwise from=120]
                child[concept color=hriSec2CompDark,text=white] { node[concept] {Action, behaviour recognition} }
                child[concept color=hriSec3,text=white] { node[concept] {Joint actions} }
            };


        \end{tikzpicture}
    }
%    \\
%    Learn -- Memorize -- Reason -- Represent -- Assess Situation
\end{frame}


\begin{frame}{Perception of human's cognitive state}

    \begin{multicols}{2}
        \resizebox{0.7\columnwidth}{!}{%
            \begin{tikzpicture}[
                    >=latex,
                every edge/.style={<-, draw, very thick}]

            \path[small mindmap, 
                level 1 concept/.append style={sibling angle=360/6}, 
                level 2 concept/.append style={sibling angle=60}, 
            concept color=hriSec1,text=white]
            node[concept] {Perception of Human's State}
                [clockwise from=60]
                child[concept color=hriSec3Dark,text=white] { node[concept]{Emotions} }
                child[concept color=hriSec2Dark,text=white] { node[concept]{Attention} }
                child[concept color=hriSec2CompDark,text=white] { node[concept] {Inference of mental models}};


        \end{tikzpicture}
    }
    \vfill
    \columnbreak

    {\Medium Emotions}

    Recognition {\tiny \Medium facial, non-verbal acoustic, verbal features / sentient analysis}\\
    Interpretation {\tiny \Medium empathy,...}

    {\Medium Attention}

    Visual attention\\
    Joint attention

    {\Medium Mental Modelling}

    User modelling\\
    (perceptual) theory of mind~\cite{lemaignan2015mutual}

    \refToContrib{Alexis' contribution}

    \end{multicols}
\end{frame}

\begin{frame}{Social Knowledge}

    \begin{multicols}{2}
        \resizebox{0.7\columnwidth}{!}{%
            \begin{tikzpicture}[
                    >=latex,
                every edge/.style={<-, draw, very thick}]

            \path[small mindmap, 
                level 1 concept/.append style={sibling angle=360/6}, 
                level 2 concept/.append style={sibling angle=60}, 
            concept color=hriSec2Comp,text=white]
            node[concept] (knowledge) {Social Knowledge} 
                [counterclockwise from=-60]
                child[concept color=hriSec1CompDark,text=white] { node[concept] {Social rules} }
                child[concept color=hriSec3Comp,text=black] { node[concept] {Social context} }
                child[concept color=hriSec3CompDark,text=white] { node[concept] {Common-sense}};


        \end{tikzpicture}
    }
    \vfill
    \columnbreak

    {\Medium Common-sense}

    Hand-coded ontologies\\
    Automated web-scrapping~\cite{tenorth2010understanding}

    {\Medium Social context}

    ?

    {\Medium Social Rules}

    Proxemics\\
    Social norms\\
    Moral norms\\
    \refToContrib{Vasanth's contribution}
    Ethics


    \end{multicols}
\end{frame}

\begin{frame}{Communication}

    \begin{multicols}{2}
        \resizebox{0.7\columnwidth}{!}{%
            \begin{tikzpicture}[
                    >=latex,
                every edge/.style={<-, draw, very thick}]

            \path[small mindmap, 
                level 1 concept/.append style={sibling angle=360/6}, 
                level 2 concept/.append style={sibling angle=60}, 
            concept color=hriSec3Comp,text=black]
            node[concept] (comm) {Communication} 
                [counterclockwise from=-30]
                child[concept color=hriSec1CompDark,text=white] { node[concept] {Verbal} }
                child[concept color=hriSec1Dark,text=white] { node[concept] {Non-verbal} };

        \end{tikzpicture}
    }
    \vfill
    \columnbreak


    Review~\cite{mavridis2015review}:
    \begin{itemize}
        \item Breaking the ``simple commands only'' barrier;
        \item Mixed initiative dialogue;
        \item Symbol grounding problem;
        \item Motor correlates
        \item ...
    \end{itemize}


    {\Medium Non-verbal}

    Behavioural alignment

    {\Medium Verbal}

    Situated dialogue~\cite{kruijff2010situated}

    \end{multicols}

\end{frame}

\begin{frame}{Interaction Dynamics}

    \begin{multicols}{2}
        \resizebox{0.7\columnwidth}{!}{%
            \begin{tikzpicture}[
                    >=latex,
                every edge/.style={<-, draw, very thick}]

            \path[small mindmap, 
                level 1 concept/.append style={sibling angle=360/6}, 
                level 2 concept/.append style={sibling angle=60}, 
            concept color=hriSec3,text=white]
            node[concept] (dynamics) {Interaction Dynamics} 
                [clockwise from=0]
                child[concept color=hriSec2Dark,text=white] { node[concept] {Long-term interaction} };

        \end{tikzpicture}
    }
    \vfill
    \columnbreak




    Behavioural alignment

    {\Medium Long-term interaction}

    Situated dialogue~\cite{kruijff2010situated}

    \end{multicols}

\end{frame}


\begin{frame}{Human-aware actions}

    \begin{multicols}{2}
        \resizebox{0.7\columnwidth}{!}{%
            \begin{tikzpicture}[
                    >=latex,
                every edge/.style={<-, draw, very thick}]

            \path[small mindmap, 
                level 1 concept/.append style={sibling angle=360/6}, 
                level 2 concept/.append style={sibling angle=60}, 
            concept color=hriSec2,text=black]
            node[concept] (action) {Human-aware actions} 
                [counterclockwise from=-30]
                child[concept color=hriSec2CompDark,text=white] { node[concept] {Action, behaviour recognition} }
                child[concept color=hriSec3,text=white] { node[concept] {Joint actions} };
        \end{tikzpicture}
    }
    \vfill
    \columnbreak




    {\Medium Action \& Behaviour recognition}

    {\Medium Joint action}


    \end{multicols}

\end{frame}

\begin{frame}{Discussion directions}
    

    \begin{itemize}
        \item Is listing functional requirements the right approach?\\ {\Medium surface competencies $\Rightarrow$ frame issue?}
        \item What are the missing ones for social interaction?\\ {\Medium imitation, self-awareness,...?}
        \item Functions interdependencies? Functional entanglement?
        \item Reference tasks? Metrics? Benchmarks? At what level are they useful? (single function? group of functions? whole architecture only?)
    \end{itemize}

    40 minutes of open discussion, plus:
    \refToContrib{Vasanth's presentation (15 min)}
    \refToContrib{Alexis' presentation (15 min)}

\end{frame}


\begin{frame}{Bibliography}
\begin{thebibliography}{10}

    \beamertemplatearticlebibitems
    \bibitem{lemaignan2015mutual}
    S. Lemaignan, P. Dillenbourg
    \newblock \doublequoted{Mutual Modelling: Inspiration for the Next Steps}
    \newblock 2015

    \beamertemplatearticlebibitems
    \bibitem{tenorth2010understanding}
    M. Tenorth, D. Nyga, M. Beetz
    \newblock \doublequoted{Understanding and Executing Instructions for Everyday Manipulation Tasks
    from the World Wide Web}
    \newblock 2010

    \beamertemplatearticlebibitems
    \bibitem{mavridis2015review}
    N. Mavridis
    \newblock \doublequoted{A review of verbal and non-verbal human–robot interactive communication}
    \newblock 2015



    \beamertemplatearticlebibitems
    \bibitem{kruijff2010situated}
    G.-J. M. Kruijff \emph{et al.}
    \newblock \doublequoted{Situated dialogue processing for human-robot interaction}
    \newblock 2010



  \end{thebibliography}
\end{frame}

\end{document}






